% Part: first-order-logic
% Chapter: proof-systems
% Section: introduction

\documentclass[../../../include/open-logic-section]{subfiles}

\begin{document}

\iftag{FOL}
      {\olfileid{fol}{prf}{int}}
      {\olfileid{pl}{prf}{int}}

\olsection{پېژندنه}

منطقونه عموماً معنٰی پوهنه او !!a{derivation} نظام دواړه لري۔ معنٰی
پوهنه د رښتياوالي، د صدق وړتيا، اعتبار او استلزام په څېر مفهومونو
سره کار لري۔ د !!{derivation} د نظامونو موخه دا ده چې د استلزام او
اعتبار د ثابتولو يوه بشپړه نحوي طريقه ورکړي۔ دا نظامونه په دې معنٰی
بشپړ نحوي دي چې په داسې يوه نظام کښې !!{derivation} يو متناهي نحوي
څيز وي، چې عموماً د !!{sentence}s يا !!{formula}s يوه لړۍ يا بل
متناهي ترتيب وي۔ د !!{derivation} ښۀ نظامونه دا خاصيت لري چې د
!!{sentence}s يا !!{formula}s د هرې ورکړل شوې لړۍ يا ترتيب
``سموالے'' په ميخانيکي توګه کتل کېدے شي۔

د لومړۍ درجې منطق تر ټولو ساده، او له تاريخي پلوه لومړني،
!!{derivation} نظامونه \emph{بديهي} وو۔ په داسې يوه نظام کښې د
!!{formula}s يوه لړۍ هله !!a{derivation} ګڼل کېږي چې په هغې کښې هر
!!{formula} يا د ``بديهي اصولو'' په يوه ټاکلي سټ کښې وي، يا د ټاکلو
``استنتاجي قاعدو'' له يوې څخه په کار اخيستو په لړۍ کښې له مخکښو
!!{formula}s څخه ترې راووځي؛ او په ميخانيکي توګه کتل کېدے شي چې آيا
!!a{formula} بديهي اصل دے او آيا د استنتاج له کومې قاعدې سره سم له
نورو !!{formula}s څخه راوتلے دے۔ د !!{derivation} بديهي نظامونه
تشريح کول اسانه دي، او له مابعدمنطقي پلوه پرې کار کول هم اسانه دي،
خو په هغو کښې !!{derivation}s لوستل او پوهېدل ګران دي او جوړول ئې هم
ګران دي۔

د !!{derivation} نور نظامونه د دې موخې دپاره جوړ شوي چې د
!!{derivation}s جوړول يا له بشپړېدو وروسته پر !!{derivation}s
پوهېدل اسانه کړي۔
بېلګې ئې طبيعي استنتاج، د رښتياوالي ونې چې د تابلو ثبوتونو په نامه
هم پېژندل کېږي، او د سېکوېنټ حساب دي۔ د !!{derivation} ځينې نظامونه
په ځانګړې توګه د ميخانيکي کولو دپاره طرح شوي دي؛ مثلاً د حل طريقه
په سافټوير کښې په اسانۍ پلې کېږي، خو پر !!{derivation}s ئې پوهېدل
تقريباً ناشوني دي۔ د !!{derivation} دا نور نظامونه زياتره
!!{derivation}s د لړيو پر ځاے د !!{formula}s د ونو په توګه ښيي۔ په
دې سره په اسانۍ ليدل کېږي چې د !!a{derivation} کومې برخې پر کومو
نورو برخو تکيه لري۔

نو د يوه ورکړل شوي منطق، لکه د لومړۍ درجې منطق، د !!{derivation}
بېل نظامونه به دا په بېلو ډولونو روښانه کړي چې !!a{sentence} کله
\emph{قضيه} ده او د نورو له ځينو څخه د !!a{sentence} د
!!{derivable} کېدو معنٰی څه ده۔ دا که په بديهي !!{derivation}s،
طبيعي استنتاجونو، سېکوېنټي !!{derivation}s، د رښتياوالي په ونو يا د
حل په ابطالونو ترسره شي، غواړو چې دا اړيکې د اعتبار او استلزام له
معنايي مفهومونو سره سمون ولري۔ د ``$!A$~قضيه ده'' دپاره
$\Proves !A$ او د ``$!A$ له~$\Gamma$ څخه !!{derivable} دے'' دپاره
``$\Gamma \Proves !A$'' ليکو۔ $\Proves$ چې هر ډول تعريف شي، غواړو
له $\Entails$ سره سمون ولري؛ يعنې:
\begin{enumerate}
\item $\Proves !A$ هله او يوازې هله چې $\Entails !A$
\item $\Gamma \Proves !A$ هله او يوازې هله چې $\Gamma \Entails !A$
\end{enumerate}
د پورته دواړو د ``يوازې هله'' لوري ته \emph{سموالے} ويل کېږي۔
!!^a{derivation} نظام هله سم دے چې !!{derivability} استلزام يا اعتبار
تضمين کړي۔ د !!{derivation} هر د کار وړ نظام بايد سم وي؛ د
!!{derivation} ناسم نظامونه هېڅ ګټه نۀ لري۔ په اصل کښې د
!!a{derivation} ټوله موخه دا ده چې د اعتبار يا استلزام نحوي تضمين
ورکړي۔ موږ به د وړاندې شويو !!{derivation} نظامونو سموالے ثابت کړو۔

د دې عکس، يعنې ``هله'' لورے، هم مهم دے: دې ته \emph{بشپړتيا} ويل
کېږي۔ د !!{derivation} بشپړ نظام دومره ځواکمن وي چې هر کله $!A$
معتبر وي، وښيي چې $!A$ قضيه ده؛ او هر کله چې
$\Gamma \Entails !A$ وي، وښيي چې $\Gamma \Proves !A$۔ بشپړتيا ثابتول
ګران دي، او ځينې منطقونه د !!{derivation} هېڅ بشپړ نظام نۀ لري۔ د
لومړۍ درجې منطق ئې لري۔ کورت ګوډل لومړے کس و چې د لومړۍ درجې منطق د
!!a{derivation} د يوه نظام بشپړتيا ئې په خپله ۱۹۲۹ز کال څېړنيزه
رساله کښې ثابته کړه۔

له !!{derivation} نظامونو سره يو بل تړلے مفهوم \emph{سازګاري} ده۔
د !!{sentence}s يو سټ هله ناسازګار بلل کېږي چې هر څه ترې
!!{derive}d کېدے شي؛ او که داسې نۀ وي، سازګار دے۔ ناسازګاري د صدق
د نۀ وړتيا نحوي سياله ده: لکه د صدق نۀ وړ سټونو غوندې، د
!!{sentence}s ناسازګار سټونه ښې نظريې نۀ جوړوي او په بنسټيز ډول
نيمګړي وي۔ د !!{sentence}s سازګار سټونه کېدے شي رښتوني يا ګټور نۀ
وي، خو لږ تر لږه د منطقي ګټورتيا له دې تر ټولو ټيټ بريد څخه تېرېږي۔
د !!{derivation} په بېلو نظامونو کښې د !!{sentence}s د سټونو د
سازګارۍ ځانګړے تعريف توپير کولے شي؛ خو لکه~$\Proves$، غواړو
سازګاري د خپل معنايي سيال، د صدق وړتيا، سره سمه وي۔ غواړو تل داسې
وي چې $\Gamma$ هله او يوازې هله سازګار وي چې د صدق وړ وي۔ دلته د
``يوازې هله'' لورے له بشپړتيا سره برابر دے، ځکه سازګاري د صدق
وړتيا تضمينوي؛ او ``هله'' لورے له سموالي سره برابر دے، ځکه د صدق
وړتيا سازګاري تضمينوي۔ په حقيقت کښې د کلاسيکي لومړۍ درجې منطق
دپاره د سموالي او بشپړتيا دواړه بڼې همارزښته دي۔

\end{document}
